\documentclass[11pt]{article}
\usepackage{fontspec}
\setmainfont{Times New Roman}
\setsansfont{Arial}
\setmonofont{Consolas}
\usepackage{amsmath,amssymb,mathtools}
\usepackage{geometry}
\usepackage{microtype}
\geometry{a4paper,margin=2.35cm}
\setlength{\parindent}{1.25em}
\setlength{\parskip}{0pt}
\makeatletter
\renewcommand\footnotesize{\@setfontsize\footnotesize{7.7}{8.7}\selectfont}
\makeatother
\emergencystretch=100em
\allowdisplaybreaks
\sloppy
\hbadness=10000
\vbadness=10000
\hfuzz=2pt
\vfuzz=10pt
\raggedbottom

\begin{document}

% Unit: Paper31 Section05 entry 2 v001. Direct Interslavic Latin translation from the German final-audited source slice, source lines 319--330 / segments P31-S0069--P31-S0070, expanded against printed scan/OCR witness pages 95--96. Footnote counter continues after Paper31 Section05 entry 1.
\setcounter{footnote}{26}

\textbf{2. \emph{Ponašanje slěda i diskriminanta pri direktnoj sumě klasov.}} Naj $\mathfrak a=\mathfrak b+\mathfrak c$. Ako za opreděljenje slěda $S_{(\mathfrak a)}(\delta)$ položimo v osnovu blokovu matricu
\[
        \begin{pmatrix}D^{\mathfrak b}&0\\0&D^{\mathfrak c}\end{pmatrix},
\]
togda, s uvagoju k punktu 1, neposrědnje čitamo: slěd elementa, vzęty vzhodno k klasu, jest suma slědov, vzętyh vzhodno k komponentnym klasam.

Zaradi izčezanja produkta $\mathfrak b\mathfrak c$ dalje dostajemo: ako $\beta$ jest element iz $\mathfrak b$, togda
\[
        S_{(\mathfrak a)}(\beta)=S_{(\mathfrak b)}(\beta);
\]
slěd $\beta$, vzęty vzhodno k klasu $(\mathfrak a)$, jest ravny slědu, vzętomu vzhodno k klasu $(\mathfrak b)$.

Iz togo, s uvagoju k §4, 4, dalje slěduje: diskriminantny ideal komponentnogo ideala $\mathfrak b$, vzęty vzhodno k klasu $(\mathfrak a)$, jest ravny diskriminantnomu idealu $\mathfrak b$, vzętomu vzhodno k komponentnomu klasu $(\mathfrak b)$.

I nakonec dostajemo: diskriminantny ideal $\mathfrak a$, vzęty vzhodno k klasu $(\mathfrak a)$, jest ravny produktu diskriminantnyh idealov komponentov, vzętyh vzhodno k komponentnym klasam. Bo ako i za tvorjenje diskriminanta, i za tvorjenje slěda položimo v osnovu bazu $\beta,\gamma$, odpovědajuču direktnoj sumě, togda baza diskriminantnogo ideala $\mathfrak a$, vzętogo vzhodno k $(\mathfrak a)$, dana jest črez
\[
\left|
\begin{array}{cc}
S_{(\mathfrak a)}(\beta_i\beta_k)&0\\
0&S_{(\mathfrak a)}(\gamma_\mu\gamma_\nu)
\end{array}
\right|
=
\left|
\begin{array}{cc}
S_{(\mathfrak b)}(\beta_i\beta_k)&0\\
0&S_{(\mathfrak c)}(\gamma_\mu\gamma_\nu)
\end{array}
\right|,
\qquad
\begin{matrix}
 i,k=1,\ldots,l,\\
 \mu,\nu=1,\ldots,m.
\end{matrix}
\]

\end{document}
